\documentclass[conference]{IEEEtran}
\usepackage{booktabs}
\usepackage{amsmath}
\usepackage{url}
\newcommand{\NAttacks}{600}
\newcommand{\NBenign}{200}
\newcommand{\TwAttackHi}{0.6\%}
\newcommand{\TwBenign}{86.0\%}
\newcommand{\CoarseBenign}{71.5\%}
\newcommand{\OutNAttacks}{200}
\newcommand{\OutTwAttack}{100.0\%}
\newcommand{\OutStrictAttack}{50.0\%}


\begin{document}

\title{Taintwatch: Span-Level Information-Flow Control\\for Tool-Using LLM Agents (Draft)}
\author{\IEEEauthorblockN{Author Name}
\IEEEauthorblockA{Affiliation\\email@example.com}}
\maketitle

\begin{abstract}
Tool-using LLM agents can be steered by instructions hidden in the data they read.
We present Taintwatch, a runtime reference monitor that labels data with integrity and
confidentiality levels, tracks those labels at character-span granularity, and checks every
tool call against a declarative policy that has a static validator.
We evaluate it on a generated benchmark of \NAttacks{} attacks and \NBenign{} benign tasks
against a mock agent that obeys injected instructions.
The default policy stops all \NAttacks{} attacks (95\% upper bound \TwAttackHi{}) while
completing \TwBenign{} of benign tasks; labeling the whole context instead of individual spans
lowers completion to \CoarseBenign{}.
Because those attacks target the sinks the policy guards, we add \OutNAttacks{} attacks that
do not: the default policy leaves \OutTwAttack{} of them through, and a stricter policy still
leaves \OutStrictAttack{}.
We discuss what this shows and what it does not.
\end{abstract}

\section{Introduction}
% TODO: motivate prompt injection against agents with tools (email, files, shell, web).
An agent that reads a web page and can also run commands mixes trusted and untrusted text in one
context. Filters that look for suspicious phrases are easy to evade and reject harmless text.
We take an information-flow view: track where each piece of data came from, and decide at the
tool boundary whether that data may reach that tool.

Contributions:
(1) span-level label tracking, which keeps trusted and untrusted parts of one string apart;
(2) a small policy language with a static validator, whose tool-coverage lint predicted a gap
that the attacks of Section~\ref{sec:eval} later exploited;
(3) a reproducible benchmark with machine-checkable success conditions, including attacks
deliberately placed outside the guarded sinks;
(4) an analysis of how many retries an attacker needs against each defense.

\section{Threat Model}
The attacker controls text that the agent reads (web pages, email, downloaded files) but not
the user's request, the policy, or the labels the runtime assigns to sources.
The attacker wants one of: running a command, leaking a secret, sending an email, writing a
file, or changing what the user is told.
We assume the labels on data are correct; Section~\ref{sec:eval} shows a case where they are not.

\section{Design}
\textbf{Labels.}
Integrity is one of untrusted, tool output, user, system; confidentiality is public, internal,
secret. Combining values lowers integrity and raises confidentiality, so taint only spreads.

\textbf{Spans.}
A tracked string stores a label per contiguous segment and keeps them through concatenation,
slicing, splitting, replacement, joining and interpolation. A call is judged by the labels of its
actual argument text, not by everything the agent has seen.

\textbf{Policy.}
Rules name a tool pattern, an optional argument, a label condition and an action
(block or confirm). A validator reports duplicate or redundant rules, impossible conditions,
and known tools that no rule covers.

\section{Benchmark}
The main suite has \NAttacks{} attacks sampled from a grid of five goals, four variants, three
carriers, six wrapper styles and three task phrasings, plus \NBenign{} benign tasks in seven
families. Success is checked on the simulated world after each run. The agent is a deterministic
mock that obeys injected directives, a worst-case model, not a measurement of real models.

\section{Evaluation}\label{sec:eval}
\textbf{Baselines.}
A phrase filter, a hand-weighted heuristic scorer, a model of spotlighting with an assumed
resistance probability, and an ablation that labels the whole context.

\begin{table}[t]
\centering
\small
\caption{Main benchmark: \NAttacks{} attacks and \NBenign{} benign tasks per defense. Attack success is a percentage with its 95\% Wilson interval.}
\label{tab:main}
\begin{tabular}{lrrr}
\toprule
Defense & Attack success [95\% CI] & Benign utility & Utility under attack \\
\midrule
none & 100.0\% [99.4, 100.0] & 100.0\% & 100.0\% \\
taintwatch & 0.0\% [0.0, 0.6] & 86.0\% & 100.0\% \\
taintwatch-coarse & 0.0\% [0.0, 0.6] & 71.5\% & 100.0\% \\
keyword & 27.0\% [23.6, 30.7] & 83.0\% & 27.0\% \\
scorer@2.5 & 31.0\% [27.4, 34.8] & 91.0\% & 31.0\% \\
spotlight@0.9 & 8.7\% [6.7, 11.2] & 100.0\% & 100.0\% \\
\bottomrule
\end{tabular}
\end{table}


\textbf{Main result.}
Table~\ref{tab:main} gives the main numbers. Taintwatch stops every attack, and so does the
whole-context ablation; the ablation differs only in utility (Table~\ref{tab:ablation}).
The phrase filter and scorer stop only attacks whose wording contains trigger phrases, which is
half of the generated wrapper styles. This depends on how the attacks were generated.

\begin{table}[t]
\centering
\small
\caption{Benign utility (\%) with span-level versus whole-context labels.}
\label{tab:ablation}
\begin{tabular}{lrr}
\toprule
Task family & Span-level & Whole-context \\
\midrule
email\_summary & 100.0\% & 0.0\% \\
file\_driven\_command & 100.0\% & 100.0\% \\
read\_file & 100.0\% & 100.0\% \\
save\_notes & 0.0\% & 0.0\% \\
summarize\_inbox & 100.0\% & 100.0\% \\
summarize\_web & 100.0\% & 100.0\% \\
user\_command & 100.0\% & 100.0\% \\
\bottomrule
\end{tabular}
\end{table}


\textbf{Attacks outside the guarded sinks.}
The main goals target sinks the default policy guards, so its result there is partly by
construction. Table~\ref{tab:outside} uses goals that avoid those sinks. The stricter policy
closes two of the four and cannot close the other two: one relies on a wrong label, and
the other changes the answer without any tool call.

\begin{table}[t]
\centering
\small
\caption{Attack success (\%) on four goals outside the default policy's guarded sinks.}
\label{tab:outside}
\begin{tabular}{lrrrr}
\toprule
Defense & answer\_poison & exfil\_get & exfil\_internal & exfil\_mislabeled \\
\midrule
none & 100.0\% & 100.0\% & 100.0\% & 100.0\% \\
taintwatch & 100.0\% & 100.0\% & 100.0\% & 100.0\% \\
taintwatch-strict & 100.0\% & 0.0\% & 0.0\% & 100.0\% \\
keyword & 42.0\% & 48.0\% & 48.0\% & 50.0\% \\
scorer@2.5 & 42.0\% & 48.0\% & 48.0\% & 50.0\% \\
\bottomrule
\end{tabular}
\end{table}


\textbf{Retries.}
Table~\ref{tab:adversary} shows the chance that an attacker trying $k$ random variants gets one
through. A defense with a modest per-variant leak rate is not safe against retries.

\begin{table}[t]
\centering
\small
\caption{Chance (\%) that $k$ random variants of the attacker's best goal include one that succeeds.}
\label{tab:adversary}
\begin{tabular}{llrrr}
\toprule
Defense & Best goal & $k{=}1$ & $k{=}10$ & $k{=}25$ \\
\midrule
none & answer\_poison & 100.0\% & 100.0\% & 100.0\% \\
taintwatch & answer\_poison & 100.0\% & 100.0\% & 100.0\% \\
taintwatch-coarse & answer\_poison & 100.0\% & 100.0\% & 100.0\% \\
taintwatch-strict & answer\_poison & 100.0\% & 100.0\% & 100.0\% \\
keyword & answer\_poison & 50.0\% & 99.9\% & 100.0\% \\
scorer@2.5 & answer\_poison & 50.0\% & 99.9\% & 100.0\% \\
spotlight@0.9 & answer\_poison & 11.6\% & 71.6\% & 96.2\% \\
\bottomrule
\end{tabular}
\end{table}


\section{Limitations}
\begin{itemize}
\item The agent is a scripted mock. Results show what the policy layer sees against a fully
compliant attacker, not how often real models are fooled. A real-model evaluation is required.
\item Attack goals and baseline thresholds were written by the authors; a different generator or
a published detector would change the baseline numbers.
\item The retry analysis is random search and understates an adaptive attacker.
\item Labels are trusted. Data that carries the wrong label defeats the monitor.
\item Utility checks verify that an action happened, not that its content was correct.
\item A confirmation rule that is auto-approved is not a defense.
\item Latency was measured on the mock and says nothing about deployment cost.
\end{itemize}

\section{Related Work}
% TODO: read and position against CaMeL and FIDES (not yet verified), prompt-injection
% benchmarks, spotlighting, and classical information-flow control.

\section{Conclusion}
Span-level labels let a monitor stop injected tool calls without rejecting tasks that merely
forward trusted values. The evidence here is from a mock agent and a generated benchmark,
and the clearest remaining work is a real-model evaluation and a policy on where data may be sent.

\end{document}
